\section{Introduction}
\label{sec:intro}
Autonomous cyber-physical systems such as cars, drones and industrial robots run a processing pipeline that typically consists of sensing, perception, planning and control. The execution path from a sensor reading to the resulting actuation must complete within an end-to-end deadline for the system to be considered functional, safe, or both. Some steps of the pipeline require hardware acceleration and routinely offload work to GPUs. Such systems are usually built on a middleware, a software layer between the application and the operating system that provides communication and composition, and thereby raises the level of abstraction at which components are written. This is what lets the hundreds of components making up the pipeline be implemented with reasonable effort.

ROS~2~\cite{ros2-standard} is one such middleware, widely used in industry and academia. It abstracts an application into three levels: the callback is the unit of computation; a node owns the interfaces through which its callbacks interact with others; and an executor dispatches the ready callbacks of the nodes assigned to it on one or more threads of an OS process.

This creates the difficulty that hinders analysis, verification and validation of such systems: three timelines run at once, and none of them sees the other two (Fig.~\ref{fig:timelines}). The OS schedules threads onto cores, each ROS~2 executor dispatches ready callbacks onto its threads, and the GPU schedules offloaded work onto its engines. In Fig.~\ref{fig:timelines}, callback 2 is a single execution to the executor, two stretches on a core to the OS, and four pieces of queued work to the GPU; the time it spends waiting---for the GPU behind another process's copy, then for a core after being preempted---appears on no single timeline. The timelines, moreover, only say where work runs; what work there is, and in what order, comes from the application's architecture.

Analysing, verifying or validating such a system therefore needs the architecture and all three timelines: which callbacks exist and how they map to nodes, executors and threads, when each ran, what GPU work it submitted, and how callbacks interact. Timing analyses of ROS~2~\cite{casini-19-reservationbased,arafat-2026-mtexecutor-cbgroups,teper-2024-end2endmtexecutoranalysisoptimization} take this information as given, and in practice it is rarely written down or current; the tools that extract it from a running system each cover a part of it, and none the GPU (Section~\ref{sec:goal}).

In this work, we present FISH\footnote{GitHub: \url{https://github.com/FatihTasyaran/FISH_INTERFERE}}\textsuperscript{,}\footnote{Documentation: \url{https://fish-interfere.readthedocs.io}}, a toolset that extracts this information automatically from a traced run of an unmodified ROS~2 application: it is transparent to the application, applies to any application, and records all three timelines on one clock, so that each can be analysed on its own and all of them together. The model is based on actual execution traces. Therefore, it does not require expert knowledge about the application that needs to be modeled, nor does it rely on potentially incomplete or outdated documentation. \textbf{Our contributions are:} \textbf{(I)}~instrumentation that discovers the deployed architecture, records every callback execution with its thread---all five kinds the executor dispatches, for rclcpp and rclpy---attributes GPU work to the callback that submitted it, and recovers how callbacks interact; \textbf{(II)}~an acquisition procedure that traces any launched application unmodified, GPU processes included, and separates its phases; \textbf{(III)}~extraction of a hierarchical model that captures each callback's GPU behaviour per invocation and can be read at any granularity, shown feasible on Autoware, an industrial-size application. Section~\ref{sec:goal} states what is extracted and why it is hard, Section~\ref{sec:method} how FISH does it, Section~\ref{sec:findings} what it yields on Autoware.

\section{Challenges and Motivation}
\label{sec:goal}
None of the three timelines can be reconstructed from its own vantage point: the OS sees threads but not what they run, an executor sees callbacks but not where they run, and the GPU sees work but not whose it is. A complete view needs all three recorded together and tied to one another, and the architecture that decides what runs and in what order. This section examines what that takes and what stands in the way.

\noindent
\textit{\textbf{Application Structure.}}
The structure is the backbone of every timeline: which components the software consists of, which functions are called when those components are triggered, and which channels they use to communicate. Knowing it completely means knowing four things. Which executor serves which nodes, and how each executor is configured: single- or multi-threaded, with how many threads, and with which callback groups, which decide which callbacks may run in parallel. The kind of every dispatchable unit of work---timer, subscription, service, client or waitable---since the executor treats them differently. Every communication pattern between callbacks: producer--consumer over a topic, request--response over a service, sample-and-hold, where a callback reads the latest value another one stored, and join, where a callback fires only once a set of inputs is complete. And the mechanism behind each: a serialised message through the middleware, a pointer handed over inside the process, or state that one callback stores and another later reads, which never becomes a message and can only be inferred.

Neither source code nor existing tools deliver this. Source code and launch files describe what may be created, not what a given run created: objects are created and destroyed while the application runs, whether a delivery stays inside the process is decided at run time, and which callback publishes on which topic is never declared, since a publisher belongs to the node and any of its callbacks may publish through it. Tools that trace the running system recover parts. ROS-Llama~\cite{blass-2021-rosllama-paper}, Abaza et al.~\cite{abaza-2024-ros2modelsynthesis} and LiME\textsuperscript{IT}~\cite{brandenburg-2026-limeit} recover callbacks and their execution times for four of the five kinds and for single-threaded executors only; CARET and its extension~\cite{kuboichi-2022-caret,peng-2022-caretcomp} record nodes, callbacks, topics and executors on top of \texttt{ros2\_tracing}~\cite{bedard-2022-ros2tracing}, whose stock tracepoints wrap neither waitables nor the deliveries inside a process. Executor configuration is recorded by none of them.

\noindent
\textit{\textbf{CPU--GPU Interplay.}}
The CPU and the GPU are two separate domains: they are traced by different tools on different clocks, and what happens on one is invisible from the other. The GPU is moreover a scheduling domain of its own, with several engines---compute and copy---each with its own queue, and a scheduler that orders the submitted work by stream and priority. Three things follow for a callback that offloads work. It is no longer atomic: what ROS~2 and the OS see as one execution becomes an alternating sequence of CPU segments and accelerator segments~\cite{enright-2024-paam}, and what a CPU-side trace records as its duration mixes computation with waiting for the GPU. Its segments are not fixed either: the GPU work a callback submits branches from invocation to invocation, so the sequence of segments---and hence the timing---differs between invocations of the same callback. And its GPU work must be attributed to it, where the only link between the two domains is the thread the GPU run-time was called from; under a multi-threaded executor that thread is decided at run time, so without the exact thread of every callback execution the callback--GPU relation cannot be established. No existing extractor records the accelerator at all: the ROS~2 tools above stop at the client library, and GPU profilers see kernels and copies but no callback.

\noindent
\textit{\textbf{Task Chains.}}
The executions spread over the three timelines are not independent: together they form task chains, from a sensor reading to an actuation, and end-to-end timing can only be computed on a chain. Yet even with the structure in hand, chains do not fall out of it. An application does not do the same thing throughout a run: some communications and executions occur only while it initialises, some only while it shuts down, and the periodic operation in between is what matters from a real-time point of view. The structure does not carry this distinction; it has to be recovered from what actually executed, and the phases separated before anything in the steady state is measured. Nor does a chain announce itself: in a deployment of hundreds of nodes the graph is entangled, a chain crosses executors, nodes and processes, and two of its links---a join, and a value sampled from stored state---leave no message and appear only as patterns in the trace. Existing tools leave the chain to the user: CARET measures a path the user names, and the extractors above produce a callback graph without the phases, the joins or the GPU segments a chain runs through.

What follows from the three challenges is the shape of the solution: one model that holds the structure, the three timelines and the relations of Fig.~\ref{fig:timelines} that tie them together, recorded on one clock and separated by phase, so that any timeline can be reconstructed alone and all of them in combination. Such a model serves more than analysis: it lets a traced run be revisited at any granularity, and it lets the same application be re-examined under a platform the trace did not run on---another core count, scheduler, executor configuration or GPU. Section~\ref{sec:method} describes how FISH builds it.
